\documentclass[11pt]{article}

\usepackage[margin=1in]{geometry}
\usepackage{amsmath, amssymb}
\usepackage{booktabs}
\usepackage{hyperref}
\usepackage{listings}
\usepackage{xcolor}
\usepackage{microtype}
\usepackage{parskip}

\hypersetup{
  colorlinks=true,
  linkcolor=blue,
  urlcolor=blue,
}

\lstset{
  basicstyle=\ttfamily\small,
  backgroundcolor=\color{gray!10},
  frame=single,
  framerule=0pt,
  breaklines=true,
  columns=fullflexible,
}

\newcommand{\ket}[1]{|#1\rangle}
\newcommand{\bra}[1]{\langle #1|}
\newcommand{\braket}[2]{\langle #1 | #2 \rangle}

\title{\textbf{superqsim}}
\author{}
\date{}

\begin{document}

\maketitle

A Python library for simulating superconducting quantum circuits. Implements three device
classes---\texttt{Transmon}, \texttt{TransmonResonator}, and
\texttt{TunableCouplerSystem}---built on \href{https://qutip.org/}{QuTiP} and NumPy. All
energies are in \textbf{GHz} with $\hbar = 1$.

\hrule
\vspace{1em}

% ============================================================
\section{Physics Background}

\subsection{The Transmon Qubit}

A transmon is a superconducting qubit based on a Josephson junction shunted by a large
capacitance. Its Hamiltonian in the Cooper-pair charge basis is:

\begin{equation}
  H = 4E_C\!\left(\hat{n} - n_g\right)^2 - E_J \cos\hat{\varphi}
\end{equation}

\begin{itemize}
  \item $\hat{n}$ --- Cooper-pair number operator (integer eigenvalues)
  \item $\hat{\varphi}$ --- superconducting phase across the junction (conjugate to $\hat{n}$)
  \item $E_C = e^2/(2C)$ --- charging energy: the electrostatic energy cost of adding one
        Cooper pair
  \item $E_J$ --- Josephson energy: the energy scale of Cooper-pair tunneling
  \item $n_g$ --- dimensionless offset charge (gate-tunable, or set by stray electric fields)
\end{itemize}

Expanding the cosine potential in the charge basis $\{\ket{n},\; n = -N,\ldots,N\}$ yields a
tridiagonal matrix:

\begin{equation}
  H_{mn} = 4E_C(m - n_g)^2\,\delta_{mn}
            \;-\; \frac{E_J}{2}\!\left(\delta_{m,n+1} + \delta_{m,n-1}\right)
\end{equation}

The diagonal gives the electrostatic energy of each charge state; the off-diagonal tunneling
terms hybridize neighboring charge states into energy bands.

\subsubsection{The $E_J/E_C$ Ratio and the Transmon Regime}

The character of the qubit is controlled by the ratio $E_J/E_C$:

\begin{itemize}
  \item \textbf{$E_J/E_C \sim 1$ (Cooper-pair box):} charge states are weakly hybridized.
        Energy levels are sensitive to $n_g$---a small stray charge shifts qubit frequency
        significantly. This ``charge noise'' is the dominant dephasing mechanism.
  \item \textbf{$E_J/E_C \gg 1$ (transmon):} deep in the cosine well, wavefunctions spread
        over many charge states. The energy bands flatten exponentially in
        $\sqrt{E_J/E_C}$, suppressing charge dispersion. Typical devices use
        $E_J/E_C \approx 50$.
\end{itemize}

The cost of this suppression is \textbf{reduced anharmonicity}: the spacing between levels
becomes more uniform. In the transmon limit the anharmonicity saturates at:

\begin{equation}
  \alpha = \omega_{12} - \omega_{01} \;\approx\; -E_C
\end{equation}

For $E_C = 0.2\,\text{GHz}$ this gives $\alpha \approx -200\,\text{MHz}$---large enough that
individual transitions can be addressed with ${\sim}10\,\text{ns}$ microwave pulses, but small
enough that the qubit frequency ($\omega_{01} \approx \text{few GHz}$) is well into the
easily-controlled microwave range.

\subsubsection{Charge Dispersion}

Charge dispersion measures how much the qubit frequency varies over a full period of $n_g$:

\begin{equation}
  \varepsilon_m = \max_{n_g} E_m - \min_{n_g} E_m
\end{equation}

It decreases exponentially with $\sqrt{E_J/E_C}$, which is why increasing $E_J/E_C$ is so
effective at reducing charge noise. The notebook's \texttt{charge\_dispersion\_sweep} demo
illustrates this directly by sweeping $n_g$ from $-1$ to $1$ and plotting the resulting energy
bands.

\hrule
\vspace{1em}

\subsection{The Transmon--Resonator System}

Coupling a transmon to a superconducting microwave resonator (a transmission-line or
lumped-element LC circuit) enables qubit control and readout via the \textbf{circuit QED}
architecture. The full Hamiltonian is:

\begin{equation}
  H = H_t \otimes \mathbb{I}_r
      + \mathbb{I}_t \otimes \omega_r a^\dagger a
      + g\,\hat{n}_t \otimes \!\left(a + a^\dagger\right)
\end{equation}

\begin{itemize}
  \item $a,\,a^\dagger$ --- resonator ladder operators
  \item $\omega_r$ --- bare resonator frequency
  \item $g$ --- capacitive coupling strength (typically $10$--$200\,\text{MHz}$)
\end{itemize}

The coupling term $g\,\hat{n}_t(a + a^\dagger)$ reflects the fact that the transmon charge
operator (electric dipole) couples to the resonator voltage (electric field).

\subsubsection{Resonant Regime ($|\Delta| \sim g$, $\Delta = \omega_{01} - \omega_r$)}

When the transmon and resonator are near resonance, they hybridize into dressed polariton
modes. The two lowest excited states repel each other, opening a gap:

\begin{equation}
  \text{vacuum-Rabi splitting} \approx 2g\,\left|\braket{0}{\hat{n}|1}\right|
\end{equation}

This avoided crossing is directly observable in the resonator frequency sweep demo.

\subsubsection{Dispersive Regime ($|\Delta| \gg g$)}

Far off resonance, the transmon and resonator remain largely separate but develop a
qubit-state-dependent resonator frequency shift. Perturbation theory gives the
\textbf{dispersive shift}:

\begin{equation}
  \chi \;\approx\; \frac{g^2\,\alpha}{\Delta(\Delta + \alpha)}
\end{equation}

where $\alpha$ is the transmon anharmonicity. The resonator frequency effectively becomes:

\begin{equation}
  \omega_r \;\to\; \omega_r + \chi\,\frac{\sigma_z}{2}
\end{equation}

so measuring the resonator transmission reveals the qubit state without directly driving the
qubit---the foundation of \textbf{dispersive readout}. The sign of $\chi$ depends on the sign
of $\Delta$ (qubit above or below the resonator).

\hrule
\vspace{1em}

\subsection{The Tunable Coupler System}

Two transmon qubits (A, B) interact via a SQUID coupler (C)---a loop of two Josephson
junctions threaded by an external magnetic flux $\Phi$. The flux tunes the effective Josephson
energy of the loop:

\begin{equation}
  E_{J,C}(\Phi) = E_{J,\max}\left|\cos\!\left(\pi\,\frac{\Phi}{\Phi_0}\right)\right|
\end{equation}

where $\Phi_0 = h/(2e)$ is the superconducting flux quantum. As $\Phi$ approaches $\Phi_0/2$,
$E_{J,C} \to 0$ and the coupler frequency drops from its maximum (at $\Phi = 0$) down toward
the qubit frequencies. This tunability is the key tool for controlling qubit--qubit coupling.

The full three-body Hamiltonian is assembled in the dressed single-qubit eigenbasis to keep the
composite Hilbert space tractable:

\begin{align}
  H &= H_A \otimes \mathbb{I}_B \otimes \mathbb{I}_C
     + \mathbb{I}_A \otimes H_B \otimes \mathbb{I}_C
     + \mathbb{I}_A \otimes \mathbb{I}_B \otimes H_C(\Phi) \notag \\
    &\quad + g_{AC}\,\hat{n}_A \otimes \mathbb{I}_B \otimes \hat{n}_C
     + g_{BC}\,\mathbb{I}_A \otimes \hat{n}_B \otimes \hat{n}_C
\end{align}

There is no direct A--B coupling term; interaction is mediated entirely through the coupler.

\subsubsection{Schrieffer--Wolff Effective Coupling}

When the coupler is far detuned from both qubits
($|\Delta_{iC}| \gg g_{iC}$, where $\Delta_{iC} = \omega_i - \omega_C$), it can be
adiabatically eliminated. The Schrieffer--Wolff (SW) perturbative transformation yields an
effective direct A--B coupling:

\begin{equation}
  g_\text{eff}(\Phi) = \frac{g_{AC}\cdot g_{BC}}{2}
                       \left(\frac{1}{\Delta_{AC}} + \frac{1}{\Delta_{BC}}\right)
\end{equation}

This formula has a striking consequence: $g_\text{eff}$ \textbf{changes sign} as $\omega_C$
sweeps through the qubit frequencies. In a real device there is typically a residual direct
capacitive coupling $g_{AB}$ between the qubits; by tuning $\Phi$ so that the coupler-mediated
term exactly cancels $g_{AB}$, the total effective coupling can be set to zero---a
``parking'' point for idle qubits. This is the operating principle of tunable-coupler
architectures used in state-of-the-art superconducting processors.

Near the degeneracy points ($\Delta_{iC} \to 0$) the perturbative approximation breaks down
and the full numerical diagonalization is required; the \texttt{flux\_sweep\_spectrum} method
provides this.

\hrule
\vspace{1em}

% ============================================================
\section{Installation}

\begin{lstlisting}[language=bash]
pip install -r requirements.txt
\end{lstlisting}

The main dependencies are \texttt{qutip}, \texttt{numpy}, \texttt{scipy}, and
\texttt{matplotlib}. The full \texttt{requirements.txt} pins all transitive dependencies for
reproducibility.

\hrule
\vspace{1em}

% ============================================================
\section{Quick Start}

\begin{lstlisting}[language=Python]
from devices import Transmon, TransmonResonator, TunableCouplerSystem
import numpy as np

# Single transmon (EJ/Ec = 50, deep transmon regime)
t = Transmon(Ec=0.2, EJ=10.0, ng=0.0, n_cutoff=15)
print(t.transition_frequency(0, 1))   # w_01 ~ 3.79 GHz
print(t.anharmonicity() * 1e3)        # alpha ~ -230 MHz

# Charge dispersion sweep
ng_vals = np.linspace(-1, 1, 300)
energies = t.charge_dispersion_sweep(ng_vals, num_levels=5)

# EJ/Ec crossover
ratios = np.linspace(1, 80, 200)
freqs, anharmon = t.ej_ec_sweep(ratios)

# Transmon coupled to a resonator (dispersive regime)
dev = TransmonResonator(Ec=0.2, EJ=10.0, omega_r=6.0, g=0.1)
print(dev.dispersive_shift * 1e3)     # chi ~ -0.43 MHz

# Two qubits with a flux-tunable coupler
sys = TunableCouplerSystem(
    Ec_A=0.20, EJ_A=10.0,
    Ec_B=0.20, EJ_B=9.5,
    Ec_C=0.30, EJ_max_C=18.0,
    g_AC=0.10, g_BC=0.10,
)
flux_vals = np.linspace(0, 0.45, 100)
g_eff = sys.effective_coupling(flux_vals)   # SW estimate vs flux
spec  = sys.flux_sweep_spectrum(flux_vals, num_levels=6)  # full numerics
\end{lstlisting}

See \texttt{transmon\_demo.ipynb} for annotated plots of all three systems.

\hrule
\vspace{1em}

% ============================================================
\section{API Reference}

\subsection{\texttt{Transmon(Ec, EJ, ng=0.0, n\_cutoff=15)}}

Single transmon qubit in the charge basis.

\begin{center}
\begin{tabular}{@{}ll@{}}
\toprule
Method / Property & Description \\
\midrule
\texttt{build\_hamiltonian()} & Tridiagonal Hamiltonian as a QuTiP \texttt{Qobj} \\
\texttt{get\_eigenspectrum(num\_levels)} & Eigenvalues ($E_0 = 0$) and eigenstates \\
\texttt{transition\_frequency(i, j)} & $E_j - E_i$; default is $\omega_{01}$ \\
\texttt{anharmonicity()} & $\alpha = \omega_{12} - \omega_{01}$ \\
\texttt{charge\_dispersion\_sweep(ng\_values, num\_levels)} & Energy bands vs $n_g$ \\
\texttt{ej\_ec\_sweep(ratio\_values)} & $\omega_{01}$ and $\alpha$ vs $E_J/E_C$ \\
\texttt{EJ\_over\_Ec} & $E_J/E_C$ ratio \\
\texttt{charge\_states} & Cooper-pair numbers $-N,\ldots,N$ \\
\bottomrule
\end{tabular}
\end{center}

\subsection{\texttt{TransmonResonator(Ec, EJ, omega\_r, g, ng=0.0, n\_cutoff=15, n\_fock=10)}}

Transmon capacitively coupled to a single-mode resonator.

\begin{center}
\begin{tabular}{@{}ll@{}}
\toprule
Method / Property & Description \\
\midrule
\texttt{build\_hamiltonian()} & Full composite Hamiltonian \\
\texttt{get\_eigenspectrum(num\_levels)} & Composite eigensystem; $E_0 = 0$ \\
\texttt{dispersive\_shift} & $\chi \approx g^2\alpha / (\Delta(\Delta + \alpha))$ \\
\texttt{resonator\_frequency\_sweep(omega\_r\_values, num\_levels)} & Spectrum vs $\omega_r$ (vacuum-Rabi crossing) \\
\texttt{transmon} & Access to the bare \texttt{Transmon} subsystem \\
\bottomrule
\end{tabular}
\end{center}

\subsection{\texttt{TunableCouplerSystem(Ec\_A, EJ\_A, \ldots, Ec\_C, EJ\_max\_C, g\_AC, g\_BC, \ldots)}}

Two transmons coupled via a flux-tunable SQUID coupler.

\begin{center}
\begin{tabular}{@{}ll@{}}
\toprule
Method / Property & Description \\
\midrule
\texttt{build\_hamiltonian(flux)} & Full 3-body Hamiltonian at dimensionless flux $\Phi/\Phi_0$ \\
\texttt{get\_eigenspectrum(flux, num\_levels)} & 3-body eigensystem; $E_0 = 0$ \\
\texttt{coupler\_EJ(flux)} & $E_{J,\max}|\cos(\pi\Phi/\Phi_0)|$ \\
\texttt{effective\_coupling(flux\_values)} & Schrieffer--Wolff $g_\text{eff}$ vs flux \\
\texttt{flux\_sweep\_spectrum(flux\_values, num\_levels)} & Full numerical spectrum vs flux \\
\bottomrule
\end{tabular}
\end{center}

\hrule
\vspace{1em}

% ============================================================
\section{Further Reading}

\begin{itemize}
  \item Koch et al., \textit{Charge-insensitive qubit design derived from the Cooper pair box},
        PRA \textbf{76}, 042319 (2007) --- original transmon proposal
  \item Blais et al., \textit{Circuit quantum electrodynamics},
        Rev.\ Mod.\ Phys.\ \textbf{93}, 025005 (2021) --- comprehensive review of circuit QED
  \item Yan et al., \textit{Tunable coupling scheme for implementing high-fidelity two-qubit
        gates}, PRA \textbf{93}, 022309 (2016) --- tunable coupler physics
\end{itemize}

\end{document}
